\subsection{Admissibility from four small weights}
\label{subsec:selmer-admissibility}

The admissibility condition in \eqref{eq:selmer-actual-admissibility}
has an elementary proof from the geometric hypotheses of the marked
model. Let $\mathcal O$ be a discrete valuation ring of characteristic
zero, with fraction field $F$, and let $K$ be an algebraically closed
extension of $F$. Let $A=\bigoplus_{j\in\mathbb Z}A_j$ be augmented
by $x:A\to\mathcal O$, where $B=A_0$ is Noetherian local,
$x(A_j)=0$ for $j\ne0$, and $A$ is generated over $B$ by finitely
many elements of weights $-2,-1,1,2$.
Let $Q$ be generated by
\[
 F_1,\ldots,F_r,F_*\in A_{-1},\quad G_j\in A_{-2},\quad
 E_k\in A_1,\quad E'_l\in A_2,\quad Z_i\in\ker(x|_B).
\]
Assume $B/(Q\cap B)$ is finite over $\mathcal O$. For $I=\ker x$,
assume the classes of the $F_i,F_*$ span $(I/I^2)_{-1}\otimes F$,
and those of the $G_j$ span $(I/I^2)_{-2}\otimes F$.

Call a point negative-only if it kills all positive weights and is
nonzero on some negative weight, and define positive-only similarly.
Assume a point of $A/Q$ killing all nonzero weights, other than $x_K$,
cannot have the same restriction to $B$ as both a negative-only and a
positive-only point. Finally, let $a\in B$ satisfy $x(a)=0$ and
$\xi(a)\ne0$ at every point $\xi:A/Q\to K$ with nonzero values of
both signs. All points here are $\mathcal O$-algebra homomorphisms.

\begin{proposition}[Admissibility]
\label{prop:selmer-admissibility}
Set $N_0=A_{-2}$, $N_1=A_2$,
$N_2=BF_*^2+A_{-2}^2A_2$, and
$\mathcal R=\bigoplus_{\boldsymbol d\ge0}
N_0^{d_0}N_1^{d_1}N_2^{d_2}t^{\boldsymbol d}$.
Let $\boldsymbol s$ consist of
$F_i^2t_0,G_jt_0,E_k^2t_1,E'_lt_1,F_*^2t_2$, where
$i=1,\ldots,r$, and let $I_{\mathrm{irr}}$ be the product of the
three coordinate-degree ideals. For every ideal $J\subset\mathcal R$
containing the $Z_i$, there is an integer $N\ge0$ such that
\[
 (\mathfrak m_B\mathcal R I_{\mathrm{irr}})^N
 \subset J+(\boldsymbol s,a).
\]
\end{proposition}

\begin{proof}
We first record an elementary consequence of the four weights:
\begin{equation}
 A_{-2n}=(A_{-2})^n,\qquad A_{2n}=(A_2)^n\quad(n\ge0).
 \label{eq:selmer-even-halves}
\end{equation}
In a monomial of even negative weight, remove opposite equal-weight
pairs. Any remaining mixed signs have weights $(-2,1)$ or $(-1,2)$;
remove their weight-zero triples. The remaining negative weight-one
factors pair up, giving the asserted factorization. The other sign is
symmetric. In particular $N_0,N_1,N_2$ are finite $B$-modules, so
$\mathcal R$ is Noetherian and standard.

Put $\mathfrak d=A_{-2}A_2\subset B$. If a character
$\sigma:B\to K$ kills $\mathfrak d$, every product of opposite
weights has zero value: for $u\in A_j,v\in A_{-j}$,
$(uv)^2\in\mathfrak d^{|j|}$. Thus $\sigma$, extended by zero in
nonzero weights, is a fixed point $\xi_0:A\to K$.
Moreover every character of $V^-=\bigoplus_{n\ge0}A_{-2n}$ over
$\sigma$ extends to a point killing all positive weights. Indeed,
for $uv\in A_{-2n}$ with $v\in A_j$, $j>0$,
\[
 (uv)^2\in (A_{-2})^{2n+j}(A_2)^j
 \subset\mathfrak d^j(A_{-2})^{2n}.
\]
The character therefore kills $V^-\cap(A_{>0})$. It descends to
the even subalgebra of $A/(A_{>0})$. The latter algebra is finite
integral over its even subalgebra, since each of its finitely many odd
generators $z$ satisfies $T^2-z^2$. Lying over and algebraic closure
extend the character to $K$. The same argument works for the positive
half, preserving any nonzero value already present.

Consider a prime $\mathfrak p$ containing $J+(\boldsymbol s,a)$
but not $I_{\mathrm{irr}}$, and write $P=\mathfrak p\cap B$.
Choose $n_i\in N_i$ with $n_it_i\notin\mathfrak p$ for each $i$.
Since $N_2\subset A_{-2}$,
\[
 (F_*^2t_0)(n_2t_2)=(F_*^2t_2)(n_2t_0)
 \quad\Longrightarrow\quad F_*^2t_0\in\mathfrak p.
\]
For a negative weight-one equation $F$ and $h\in A_1$, the identity
$(hF)^2(n_0t_0)=(F^2t_0)(h^2n_0)$ gives $hF\in P$.
For $G\in A_{-2}$ and $h\in A_2$, use
$(hG)(n_0t_0)=(Gt_0)(hn_0)$; use $n_1t_1$ for the positive
equations. Together with the $Z_i$, these identities kill every
weight-zero term of the ideal $Q$, so $Q\cap B\subset P$.

Consequently $B/P$ is finite over $\mathcal O$. If $P$ meets
$\mathcal O\setminus\{0\}$, then $B/P$ is a finite-dimensional
domain over the residue field, hence a field: $P=\mathfrak m_B$.
Otherwise its fraction field embeds in $K$, giving
$\sigma:B\to K$ with kernel $P$. There is also a character
$\theta:\mathcal R\to K$ over $\sigma$, killing $\mathfrak p$
and nonzero on all $n_it_i$. To see this, localize the domain
$\mathcal R/\mathfrak p$ over $\operatorname{Frac}(B/P)$, extend
scalars to $K$, and invert $\prod_i n_it_i$. This is a nonzero
finite-type $K$-algebra: field extension preserves the injectivity of
multiplication by every power of that product. The weak Nullstellensatz
supplies the desired character.

If $\sigma(\mathfrak d)\ne0$, the graded algebra
$(A/Q)\otimes_{B,\sigma}K$ has degree-zero part $K$ and is a
nonzero finite-type $K$-algebra. A $K$-point of it is two-sided because
some product from $A_{-2}A_2$ is nonzero. It has $a$-value zero,
contrary to the hypothesis on $a$.

If $\sigma(\mathfrak d)=0$, restrict $\theta$ to the first and
second axes of $\mathcal R$. By \eqref{eq:selmer-even-halves} these
are the two even half-algebras. The extensions just constructed give
nontrivial negative-only and positive-only points $\xi^-,\xi^+$.
They kill $Q$: the negative equations follow from
$F_i^2t_0,F_*^2t_0,G_jt_0\in\mathfrak p$, and the positive
ones from the corresponding second-axis sections. All three points
$\xi^-,\xi^+,\xi_0$ restrict to $\sigma$, and $\xi_0$ also kills
$Q$. If $\xi_0\ne x_K$, fixed-point separation is contradicted.

If $\xi_0=x_K$, the negative-only point kills every nonnegative
weight in $I$. Over $F$, each weight $-1$ product in $I^2$ has a
nonnegative-weight factor. Cotangent spanning and the vanishing of the
$F_i,F_*$ force $\xi^-(A_{-1})=0$. A weight $-2$ product in $I^2$
has either a nonnegative-weight factor or two weight $-1$ factors.
Spanning by the $G_j$ therefore forces $\xi^-(A_{-2})=0$, contrary
to its nonzero value on $n_0$.

Every prime over $J+(\boldsymbol s,a)$ thus contains either
$I_{\mathrm{irr}}$ or $\mathfrak m_B\mathcal R$. Their product
lies in the radical, and Noetherianity gives the required power.
\end{proof}

This proves \eqref{eq:selmer-actual-admissibility} for the chart ideal,
which contains the degree-zero equations. Together with the preceding
Koszul and leading-term calculations it supplies the lower comparison
under the stated model hypotheses. Constructing that arithmetic model
and proving its point conditions remain separate inputs.
